\documentclass[../../数学分析培优讲义.tex]{subfiles}

\begin{document}

\setcounter{chapter}{4}
\pagestyle{fancy}

\chapter{不定积分}


\section{不定积分计算方法}


\subsection{基本方法}

\begin{theorem}[\markstar\ 常见三角函数公式]
\begin{enumerate}[label=\arabic*.]
    \item 积化和差公式
    \begin{align*}
        \sin x\cos y&=\dfrac{1}{2}[\sin(x+y)+\sin(x-y)],
        &\cos x\sin y&=\dfrac{1}{2}[\sin(x+y)-\sin(x-y)],\\
        \cos x\cos y&=\dfrac{1}{2}[\cos(x+y)+\cos(x-y)],
        &\sin x\sin y&=\dfrac{1}{2}[\cos(x-y)-\cos(x+y)].
    \end{align*}
    \item 和差化积公式
    \begin{align*}
        \sin\alpha+\sin\beta
        &=2\sin\dfrac{\alpha+\beta}{2}\cos\dfrac{\alpha-\beta}{2},\\
        \sin\alpha-\sin\beta
        &=2\cos\dfrac{\alpha+\beta}{2}\sin\dfrac{\alpha-\beta}{2},\\
        \cos\alpha+\cos\beta
        &=2\cos\dfrac{\alpha+\beta}{2}\cos\dfrac{\alpha-\beta}{2},\\
        \cos\alpha-\cos\beta
        &=-2\sin\dfrac{\alpha+\beta}{2}\sin\dfrac{\alpha-\beta}{2}.
    \end{align*}
    \item 万能公式
    \[
        \tan\alpha=\dfrac{2\tan\dfrac{\alpha}{2}}{1-\tan^2\dfrac{\alpha}{2}},
        \qquad
        \sin\alpha=\dfrac{2\tan\dfrac{\alpha}{2}}{1+\tan^2\dfrac{\alpha}{2}},
        \qquad
        \cos\alpha=\dfrac{1-\tan^2\dfrac{\alpha}{2}}{1+\tan^2\dfrac{\alpha}{2}}.
    \]
    \item 三角恒等式
    \[
        \sec^2x=1+\tan^2x,
        \qquad
        \csc^2x=1+\cot^2x.
    \]
    \[
        \cos\dfrac{\alpha}{2}\cos\dfrac{\alpha}{2^2}\cdots
        \cos\dfrac{\alpha}{2^n}
        =\dfrac{\sin\alpha}{2^n\sin\dfrac{\alpha}{2^n}},
        \qquad
        \sum\limits_{k=1}^n\sin kx
        =\dfrac{\cos\dfrac{\pi}{2}-\cos\left(n+\dfrac{1}{2}\right)x}
        {2\sin\dfrac{x}{2}}.
    \]
    \[
        \sum\limits_{k=1}^n\cos kx
        =\dfrac{\cos\left(n+\dfrac{1}{2}\right)x-\sin\dfrac{x}{2}}
        {2\sin\dfrac{x}{2}}.
    \]
    \item 求导公式
    \[
        (\sec x)'=\tan x\sec x,
        \qquad
        (\tan x)'=\sec^2x.
    \]
\end{enumerate}
\end{theorem}

\begin{theorem}[\markstar\ 部分初等函数积分公式]
\begin{enumerate}[label=(\arabic*)]
    \item $\displaystyle\int x^\alpha\,\mathrm{d}x=\dfrac{x^{\alpha+1}}{\alpha+1}+C\quad(\alpha\neq-1)$;
    \item $\displaystyle\int\dfrac{1}{x}\,\mathrm{d}x=\ln|x|+C$;
    \item $\displaystyle\int e^x\,\mathrm{d}x=e^x+C$;
    \item $\displaystyle\int a^x\,\mathrm{d}x=\dfrac{a^x}{\ln a}+C\quad(a>0,a\neq1)$;
    \item $\displaystyle\int\cos x\,\mathrm{d}x=\sin x+C$;
    \item $\displaystyle\int\sin x\,\mathrm{d}x=-\cos x+C$;
    \item $\displaystyle\int\tan x\,\mathrm{d}x=-\ln|\cos x|+C$;
    \item $\displaystyle\int\cot x\,\mathrm{d}x=\ln|\sin x|+C$;
    \item $\displaystyle\int\sec^2x\,\mathrm{d}x=\tan x+C$;
    \item $\displaystyle\int\csc^2x\,\mathrm{d}x=-\cot x+C$;
    \item $\displaystyle\int\dfrac{\mathrm{d}x}{\sqrt{a^2-x^2}}=\arcsin\dfrac{x}{a}+C$;
    \item $\displaystyle\int\dfrac{\mathrm{d}x}{x^2+a^2}=\dfrac{1}{a}\arctan\dfrac{x}{a}+C$;
    \item $\displaystyle\int\dfrac{\mathrm{d}x}{x^2-a^2}=\dfrac{1}{2a}\ln\left|\dfrac{x-a}{x+a}\right|+C$;
    \item $\displaystyle\int\dfrac{\mathrm{d}x}{\sqrt{x^2\pm a^2}}=\ln\left|x+\sqrt{x^2\pm a^2}\right|+C$;
    \item $\displaystyle\int\sqrt{a^2-x^2}\,\mathrm{d}x
    =\dfrac{1}{2}\left(x\sqrt{a^2-x^2}+a^2\arcsin\dfrac{x}{a}\right)+C$;
    \item $\displaystyle\int\sqrt{x^2+a^2}\,\mathrm{d}x
    =\dfrac{1}{2}\left[x\sqrt{x^2+a^2}+a^2\ln\left(x+\sqrt{x^2+a^2}\right)\right]+C$;
    \item $\displaystyle\int e^{ax}\cos bx\,\mathrm{d}x
    =\dfrac{e^{ax}}{a^2+b^2}(a\cos bx+b\sin bx)$;
    \item $\displaystyle\int e^{ax}\sin bx\,\mathrm{d}x
    =\dfrac{e^{ax}}{a^2+b^2}(a\sin bx-b\cos bx)$.
\end{enumerate}
\end{theorem}

\begin{theorem}[第一积分换元法]
若 $f(u)$ 在区间 $J$ 上有原函数 $F(u)$,即
\[
    \int f(u)\,\mathrm{d}u=F(u)+C,
    \qquad u\in J,
\]
则 $F(\varphi(x))$ 是 $f(\varphi(x))\varphi'(x)$ 在区间 $I$ 上的原函数:
\[
    \int f(\varphi(x))\varphi'(x)\,\mathrm{d}x
    =F(\varphi(x))+C,
    \qquad x\in I,
\]
其中 $u=\varphi(x)$ 是 $I$ 上的可微函数,且 $\varphi(I)\subseteq J$.
\end{theorem}

\begin{remark}
关键在于将被积函数凑成 $f[\varphi(x)]\varphi'(x)\,\mathrm{d}x$ 的形式,故俗称凑微分法.
\end{remark}

\begin{corollary}[\markstar]
设 $f(x)$ 在区间 $I$ 上可微,$f(x)\neq0$($x\in I$),则:
\begin{enumerate}[label=(\arabic*)]
    \item $\displaystyle\int\dfrac{f'(x)}{f(x)}\,\mathrm{d}x=\ln|f(x)|+C$;
    \item 若 $f(x)>0$($x\in I$),则
    \[
        \int\dfrac{f'(x)}{\sqrt{f(x)}}\,\mathrm{d}x=2\sqrt{f(x)}+C,
        \qquad
        \int f'(x)\sqrt{f(x)}\,\mathrm{d}x
        =\dfrac{2}{3}f^{\frac{3}{2}}(x)+C.
    \]
\end{enumerate}
\end{corollary}

\begin{theorem}[第二积分换元法]
设 $G(t)$ 是 $f(\varphi(t))\varphi'(t)$ 在区间 $I$ 的原函数,
\[
    \int f(\varphi(t))\varphi'(t)\,\mathrm{d}t=G(t)+C.
\]
则 $G(\varphi^{-1}(x))$ 是 $f(x)$ 在区间 $J$ 上的原函数:
\[
    \int f(x)\,\mathrm{d}x=G(\varphi^{-1}(x))+C,
    \qquad x\in J.
\]
\end{theorem}

\begin{example}
求下列不定积分:
\begin{enumerate}[label=(\arabic*)]
    \item $\displaystyle I=\int\dfrac{\mathrm{d}x}{3x^2+3x-1}$;
    \item $\displaystyle I=\int\dfrac{\mathrm{d}x}{\sqrt{3x^2+3x-1}}$;
    \item $\displaystyle I=\int\dfrac{\mathrm{d}x}{x^2-5x+7}$.
\end{enumerate}
\end{example}
\exampleblank

\begin{example}
求下列不定积分:
\begin{enumerate}[label=(\arabic*)]
    \item $\displaystyle I=\int\dfrac{a^x-a^{-x}}{a^x+a^{-x}}\,\mathrm{d}x\quad(a>0)$;
    \item $\displaystyle I=\int\dfrac{\ln x+1}{1+x\ln x}\,\mathrm{d}x$;
    \item $\displaystyle I=\int\dfrac{\sin2x}{a^2\sin^2x+b^2\cos^2x}\,\mathrm{d}x\quad(a^2\neq b^2)$;
    \item $\displaystyle I=\int(3x-1)\sqrt{3x^2-2x+7}\,\mathrm{d}x$.
\end{enumerate}
\end{example}
\exampleblank

\begin{example}
求下列不定积分:
\begin{enumerate}[label=(\arabic*)]
    \item $\displaystyle I=\int(\tan x-\sec x)^2\,\mathrm{d}x$;
    \item $\displaystyle I=\int\sin^4x\,\mathrm{d}x$;
    \item $\displaystyle I=\int\dfrac{\sin^4x}{\cos^2x}\,\mathrm{d}x$;
    \item $\displaystyle I=\int\sin mx\cos nx\,\mathrm{d}x$.
\end{enumerate}
\end{example}
\exampleblank

\begin{example}
求下列不定积分:
\begin{enumerate}[label=(\arabic*)]
    \item $\displaystyle I=\int\dfrac{\mathrm{d}x}{x^4+x}$;
    \item $\displaystyle I=\int\dfrac{x\,\mathrm{d}x}{\sqrt{1-4x^2}}$;
    \item $\displaystyle I=\int\dfrac{\mathrm{d}x}{x^4\sqrt{1+x^2}}$;
    \item $\displaystyle I=\int\dfrac{\mathrm{d}x}{x\sqrt{x^2-1}}$.
\end{enumerate}
\end{example}
\exampleblank

\begin{example}
求下列函数的不定积分:
\begin{enumerate}[label=(\arabic*)]
    \item $\displaystyle I=\int\dfrac{\mathrm{d}x}{x\ln x}$;
    \item $\displaystyle I=\int\dfrac{\mathrm{d}x}{x^8(1+x^2)}$.
\end{enumerate}
\end{example}
\exampleblank

\begin{example}
求下列不定积分:
\begin{enumerate}[label=(\arabic*)]
    \item $\displaystyle I=\int\dfrac{\mathrm{d}x}{\sin x}$;
    \item $\displaystyle I=\int\dfrac{\mathrm{d}x}{\cos x}$;
    \item $\displaystyle I=\int\dfrac{\cos x}{\sqrt{\cos2x}}\,\mathrm{d}x$;
    \item $\displaystyle I=\int\dfrac{\mathrm{d}x}{3\cos^2x+4\sin^2x}$;
    \item $\displaystyle I=\int\dfrac{\sin x}{1+\sin x}\,\mathrm{d}x$;
    \item $\displaystyle I=\int\dfrac{\mathrm{d}x}{2+\cos^2x}$.
\end{enumerate}
\end{example}
\exampleblank

\begin{theorem}[\marktwostar\ 分部积分法]
设 $u(x),v(x)$ 在区间 $I$ 上可微,若 $v(x),u'(x)$ 在 $I$ 上有原函数,则 $u(x),v'(x)$ 在 $I$ 上也有原函数,而且
\[
    \int u(x)v'(x)\,\mathrm{d}x
    =u(x)v(x)-\int v(x)u'(x)\,\mathrm{d}x.
\]
\end{theorem}

\begin{remark}
即将 $uv'$ 原函数求解问题转化为更易求解的 $vu'$ 原函数问题.而 $u,v$ 究竟如何选择呢?优先选 $u$ 的一般顺序为(反对幂三指):
\[
    \text{反三角函数}\longrightarrow\text{对数函数}\longrightarrow\text{幂函数}
    \longrightarrow\text{三角函数}\longrightarrow\text{指数函数}.
\]
\end{remark}

\begin{example}
求下列不定积分:
\begin{enumerate}[label=(\arabic*)]
    \item $\displaystyle I=\int\ln x\,\mathrm{d}x$;
    \item $\displaystyle I=\int xe^x\,\mathrm{d}x$;
    \item $\displaystyle I=\int\ln(1+\sqrt{x})\,\mathrm{d}x$;
    \item $\displaystyle I=\int\left(\dfrac{\ln x}{x}\right)^2\,\mathrm{d}x$.
\end{enumerate}
\end{example}
\exampleblank

\begin{example}
求下列函数的不定积分:
\begin{enumerate}[label=(\arabic*)]
    \item $\displaystyle I=\int x\cos x\,\mathrm{d}x$;
    \item $\displaystyle I=\int(2x+3x^2)\arctan x\,\mathrm{d}x$;
    \item $\displaystyle I=\int\dfrac{x\arcsin x}{\sqrt{1-x^2}}\,\mathrm{d}x$;
    \item $\displaystyle I=\int\dfrac{\ln\sin x}{\sin^2x}\,\mathrm{d}x$;
    \item $\displaystyle I=\int\sin(\ln x)\,\mathrm{d}x$.
\end{enumerate}
\end{example}
\exampleblank



\subsection{特殊方法}


\methodsection{(A) 配对法}

当要求不定积分 $I(x)$,但 $I(x)$ 常规办法不易求得时,对偶地考虑与其相关的 $J(x)$,使得 $aI(x)+bJ(x)$ 与 $cI(x)+dJ(x)$ 更易求得,即可由此解出 $I(x)$.

\begin{proposition}[\markstar]
对于形如
\[
    I=\int\dfrac{a\sin x+b\cos x}{c\sin x+d\cos x}\,\mathrm{d}x
\]
的不定积分,取 $A,B$ 使得
\[
    A(c\sin x+d\cos x)+B(c\sin x+d\cos x)
    =a\sin x+b\cos x.
\]
则有
\[
    I=Ax+B\ln|c\sin x+d\cos x|+C.
\]
更进一步,类似上述方法可解所有形如
\[
    \int\dfrac{p\sin x+q\cos x+s}{a\sin x+b\cos x+c}\,\mathrm{d}x
\]
的不定积分.
\end{proposition}

\begin{example}
求下列不定积分:
\begin{enumerate}[label=(\arabic*)]
    \item $\displaystyle\int\dfrac{\sin x}{\sin x+\cos x}\,\mathrm{d}x$;
    \item $\displaystyle\int\dfrac{\mathrm{d}x}{3+5\tan x}$;
    \item $\displaystyle\int\dfrac{\mathrm{d}x}{\sin x+2\cos x+3}$.
\end{enumerate}
\end{example}
\exampleblank

\begin{example}
利用配对法求下列不定积分:
\begin{enumerate}[label=(\arabic*)]
    \item $\displaystyle\int\dfrac{\mathrm{d}x}{1+x^4}$;
    \item $\displaystyle\int\dfrac{\mathrm{d}x}{1+x^2+x^4}$;
    \item $\displaystyle\int\dfrac{\mathrm{d}x}{1+x^3}$.
\end{enumerate}
\end{example}
\exampleblank



\methodsection{(B) 递推法}

为了计算 $\displaystyle I_n=\int f^n(x)\,\mathrm{d}x$,可以用各种方法将它化成参数值较小的 $\displaystyle I_{n-k}=\int f^{n-k}(x)\,\mathrm{d}x$($0<k\leq n$),并且如此进行下去,直到最后把问题化为参数值最小的一个或几个积分.

\begin{proposition}[\markstar]
设
\[
    I(m,n)=\int\cos^m x\sin^n x\,\mathrm{d}x,
\]
则
\[
    I(m,n)
    =\dfrac{\cos^{m-1}x\sin^{n+1}x}{m+n}
    +\dfrac{n-1}{m+n}I(m-2,n).
\]
\end{proposition}

\begin{example}
求下列不定积分的递推公式:
\begin{enumerate}[label=(\arabic*)]
    \item $\displaystyle I_n=\int\tan^n x\,\mathrm{d}x$;
    \item $\displaystyle I_n=\int\sec^n x\,\mathrm{d}x$;
    \item $\displaystyle I_n=\int\dfrac{\sin nx}{\sin x}\,\mathrm{d}x$.
\end{enumerate}
\end{example}
\exampleblank

\begin{example}
求下列不定积分的递推公式:
\begin{enumerate}[label=(\arabic*)]
    \item $\displaystyle I_n=\int\dfrac{\mathrm{d}x}{(1+x^2)^n}$;
    \item $\displaystyle I_n=\int\ln^n x\,\mathrm{d}x$;
    \item $\displaystyle I_{m,n}=\int x^m\ln^n x\,\mathrm{d}x$.
\end{enumerate}
\end{example}
\exampleblank


\newpage

\section{几类可积函数}


\subsection{有理分式}

对于有理真分式,总可分解为形如下列四种最简真分式的组合:
\[
    \dfrac{A}{x-a},\qquad
    \dfrac{A}{(x-a)^k}\quad(k\geq2),\qquad
    \dfrac{Ax+B}{x^2+px+q},\qquad
    \dfrac{Ax+B}{(x^2+px+q)^k}
\]
其中 $p^2-4q<0,k\geq2$.

我们有:
\begin{enumerate}[label=(\arabic*)]
    \item
    \[
        \int\dfrac{A}{x-a}\,\mathrm{d}x=A\ln|x-a|+C.
    \]
    \item
    \[
        \int\dfrac{A}{(x-a)^k}\,\mathrm{d}x
        =\dfrac{A(x-a)^{1-k}}{1-k}+C.
    \]
    \item
    \begin{align*}
        \int\dfrac{Ax+B}{x^2+px+q}\,\mathrm{d}x
        &=\dfrac{A}{2}\int\dfrac{2x+p}{x^2+px+q}\,\mathrm{d}x
        +\left(B-\dfrac{Ap}{2}\right)\int\dfrac{\mathrm{d}x}{x^2+px+q}\\
        &=\dfrac{A}{2}\ln|x^2+px+q|
        +\left(B-\dfrac{Ap}{2}\right)
        \int\dfrac{\mathrm{d}x}{\left(x+\dfrac{p}{2}\right)^2+q-\dfrac{p^2}{4}}\\
        &=\dfrac{A}{2}\ln|x^2+px+q|
        +\dfrac{2B-Ap}{\sqrt{4q-p^2}}
        \arctan\dfrac{2x+p}{\sqrt{4q-p^2}}+C.
    \end{align*}
    \item 可得递推公式(略).
\end{enumerate}

\begin{example}
求下列不定积分:
\begin{enumerate}[label=(\arabic*)]
    \item $\displaystyle\int\dfrac{\mathrm{d}x}{x(x^3+1)^2}$;
    \item $\displaystyle\int\dfrac{x(x^2+3)}{(x^2-1)(x^2+1)^2}\,\mathrm{d}x$;
    \item $\displaystyle\int\dfrac{2x^4+5x^2-2}{2x^3-x-1}\,\mathrm{d}x$;
    \item $\displaystyle\int\dfrac{\mathrm{d}x}{\cos^7x}$.
\end{enumerate}
\end{example}
\exampleblank



\subsection{无理函数}

\begin{proposition}[\markstar]
对于下列四种类型的不定积分,有如下处理技巧:
\begin{enumerate}[label=\textcircled{\arabic*}]
    \item $R\left(x,x^{\frac{m}{n}},\ldots,x^{\frac{r}{s}}\right)$ 型:记 $\dfrac{m}{n},\ldots,\dfrac{r}{s}$ 的公分母为 $k$,令 $x=t^k$.
    \item $R\left[x,\left(\dfrac{ax+b}{cx+d}\right)^{\frac{m}{n}},\ldots,
    \left(\dfrac{ax+b}{cx+d}\right)^{\frac{t}{s}}\right]$ 型:类似 \textcircled{1},记各分数指数的公分母为 $k$,令
    \[
        \dfrac{ax+b}{cx+d}=t^k.
    \]
    \item $R(x,\sqrt{ax^2+bx+c})$ 型:根据 $a,b,c$ 的系数,考虑以下三种 Euler 变换:

    Euler 第一变换:若 $a>0$,则令
    \[
        \sqrt{ax^2+bx+c}=\pm\sqrt{a}\,x+t.
    \]

    Euler 第二变换:若 $c>0$,则令
    \[
        \sqrt{ax^2+bx+c}=xt+\sqrt{c}.
    \]

    Euler 第三变换:若 $ax^2+bx+c$ 有两个实根,即 $ax^2+bx+c=a(x-\alpha)(x-\beta)$,则令
    \[
        \sqrt{ax^2+bx+c}=t(x-\alpha).
    \]
    \item $x^m(ax^n+b)^p$ 型:令 $x^n=t$,化为前述类型之一.
\end{enumerate}
\end{proposition}

\begin{remark}
事实上,\textcircled{3} 中利用 Euler 第一、三变换足矣,因为若 $\Delta\geq0$ 可用第三变换;若 $\Delta<0$,由
\[
    ax^2+bx+c=\dfrac{1}{4a}\left[(2ax+b)^2+(4ac-b^2)\right]
\]
可得 $ax^2+bx+c$ 与 $a$ 符号相同,故可用第一变换.
\end{remark}

\begin{example}
求下列不定积分:
\begin{enumerate}[label=(\arabic*)]
    \item $\displaystyle\int\dfrac{x+\sqrt[3]{x^2}+6\sqrt{x}}{x(1+\sqrt[3]{x})}\,\mathrm{d}x$;
    \item $\displaystyle\int\dfrac{\mathrm{d}x}{\sqrt{x+1}+\sqrt[3]{x+1}}$;
    \item $\displaystyle\int\sqrt[3]{\dfrac{2-x}{2+x}}\dfrac{\mathrm{d}x}{(2-x)^2}$;
    \item $\displaystyle\int\dfrac{1}{x}\sqrt{\dfrac{x+2}{x-2}}\,\mathrm{d}x$.
\end{enumerate}
\end{example}
\exampleblank

\begin{example}
求下列不定积分:
\begin{enumerate}[label=(\arabic*)]
    \item $\displaystyle\int\dfrac{\sqrt{x^2+2x+3}}{x}\,\mathrm{d}x$;
    \item $\displaystyle\int\dfrac{\mathrm{d}x}{x^2\sqrt{4x^2-3x-1}}$;
    \item $\displaystyle\int\dfrac{1-\sqrt{1+x+x^2}}{x\sqrt{1+x+x^2}}\,\mathrm{d}x$;
    \item $\displaystyle\int\dfrac{\mathrm{d}x}{\sqrt{x^2+3x-4}}$.
\end{enumerate}
\end{example}
\exampleblank

\begin{example}
求下列不定积分:
\begin{enumerate}[label=(\arabic*)]
    \item $\displaystyle\int x^5(a^3+x^3)^{\frac{1}{3}}\,\mathrm{d}x$;
    \item $\displaystyle\int\dfrac{(x^6+a^6)^3}{x^{11}}\,\mathrm{d}x$.
\end{enumerate}
\end{example}
\exampleblank



\subsection{三角函数有理式}

\begin{proposition}[\markstar]
由万能公式,令 $t=\tan\dfrac{x}{2}$ 可得
\[
    \int R(\cos x,\sin x)\,\mathrm{d}x
    =\int R\left(\dfrac{1-t^2}{1+t^2},\dfrac{2t}{1+t^2}\right)
    \dfrac{2}{1+t^2}\,\mathrm{d}t.
\]
仍为有理函数,由此可转化为有理函数的不定积分.但实际计算中,该方法不一定是最简便的.因此,遇到下列情形,应灵活设计变量替换:
\begin{enumerate}[label=\textcircled{\arabic*}]
    \item 若有 $R(\cos x,-\sin x)=-R(\cos x,\sin x)$,则可令 $t=\cos x$;
    \item 若有 $R(-\cos x,\sin x)=-R(\cos x,\sin x)$,则可令 $t=\sin x$;
    \item 若有 $R(-\cos x,-\sin x)=-R(\cos x,\sin x)$,则可令 $t=\tan x$.
\end{enumerate}
\end{proposition}

\begin{example}
求下列不定积分:
\begin{enumerate}[label=(\arabic*)]
    \item $\displaystyle\int\dfrac{\mathrm{d}x}{a+b\tan x}$;
    \item $\displaystyle\int\dfrac{\mathrm{d}x}{\sin x\cos2x}$;
    \item $\displaystyle\int\dfrac{\cos2x}{\sin^4x+\cos^4x}\,\mathrm{d}x$;
    \item $\displaystyle\int\dfrac{\mathrm{d}x}{1-\sin x}$;
    \item $\displaystyle\int\dfrac{\sin x}{1+\sin x}\,\mathrm{d}x$;
    \item $\displaystyle\int\dfrac{\mathrm{d}x}{5+3\sin x+4\cos x}$;
    \item $\displaystyle\int\dfrac{\sin x}{1+\sin x+\cos x}\,\mathrm{d}x$.
\end{enumerate}
\end{example}
\exampleblank

\end{document}
